\documentclass{article}
\usepackage{arxiv}

\usepackage[utf8]{inputenc}
\usepackage[T1]{fontenc}
\usepackage{hyperref}
\usepackage{url}
\usepackage{booktabs}
\usepackage{amsmath, amssymb}
\usepackage{bm}
\usepackage{natbib}
\usepackage{graphicx}

% The emblem lives in the repository-wide assets directory, two levels up.
\graphicspath{{../../assets/}{../../assets/logos/}{../../assets/biota/}{../../assets/diagrams/}}

\title{The Mathematics and Design of FlameSkimmer}

% arxiv.sty labels every paper "A Preprint" in the running head and under the
% title. This is an internal technical report, not a preprint of a submission.
\renewcommand{\headeright}{Technical Report}
\renewcommand{\undertitle}{Technical Report}

% arxiv.sty's title block sets a 24pt strut above the author names and 0.4in
% below them, which leaves the byline stranded between "Technical Report" and
% the date. This is the same layout with those two gaps halved. Overridden here
% rather than edited in arxiv.sty, which is a vendored file.
\makeatletter
\renewcommand{\@maketitle}{%
  \vbox{%
    \hsize\textwidth
    \linewidth\hsize
    \vskip 0.1in
    \@toptitlebar
    \centering
    {\LARGE\sc \@title\par}
    \@bottomtitlebar
    \textsc{\undertitle}\\
    \vskip 0.1in
    \def\And{%
      \end{tabular}\hfil\linebreak[0]\hfil%
      \begin{tabular}[t]{c}\bf\rule{\z@}{12\p@}\ignorespaces%
    }
    \def\AND{%
      \end{tabular}\hfil\linebreak[4]\hfil%
      \begin{tabular}[t]{c}\bf\rule{\z@}{12\p@}\ignorespaces%
    }
    \begin{tabular}[t]{c}\bf\rule{\z@}{12\p@}\@author\end{tabular}%
  \vskip 0.2in \@minus 0.1in \center{\@date}   \vskip 0.2in
  }
}
\makeatother

\author{
  Daniel Boupha Christensen \\
  Environmental, Health, and Safety, University of California, Berkeley \\
  The Green Initiative Fund \\
}

\begin{document}

% Masthead. Licence at the left, group emblem at the right. This is an internal
% report, so the block carries the licence and the project mark only, and
% deliberately no journal, volume, page range or DOI furniture.
\noindent
\begin{minipage}[c]{0.74\textwidth}
  \footnotesize
  Copyright 2024--2026 Strawberry Creek Monitoring Group contributors.\\
  Licensed under the Apache License, Version 2.0.\\
  \url{https://www.apache.org/licenses/LICENSE-2.0}
\end{minipage}%
\hfill
\begin{minipage}[c]{0.16\textwidth}
  \raggedleft
  \includegraphics[width=\linewidth]{SCMGemblem.png}
\end{minipage}

\maketitle

\begin{abstract}
FlameSkimmer is the third detector built for Strawberry Creek, and its job is
to know when it is out of its depth. DuskCrayfish gives a point forecast, and
CobbleShoal adds a scale that says how noisy a sensor usually is, but neither
can say that it has never seen conditions like the ones in front of it. A storm,
a change of season or a long outage is forced into the same yes or no as an
ordinary afternoon. FlameSkimmer keeps CobbleShoal's encoder unchanged and
replaces its forecast head with a small Bayesian neural network, one that
learns a distribution over its weights rather than a single value for each.
Every forecast then carries two variances: the noise the sensor is expected to
show, and the disagreement between plausible versions of the model. This
document sets out the mathematics of that head, how it is trained, how its two
variances feed the detection channels, and what the first trained checkpoint
measured. The short version of the last part is that the machinery works and
the uncertainty it reports does not yet respond to storms, for a reason that
is now understood and can be fixed.
\end{abstract}

\section{Introduction}
\label{sec:introduction}

A forecasting detector flags a reading when it lands somewhere the model did not
expect. That works when the model's expectations are good. It goes wrong when
the creek does something the model has rarely or never seen, because the model
still produces a forecast, still produces a scale around it, and reports both
with the same confidence it would have on a dry Tuesday in August. The
resulting score cannot tell an unusual reading from an unusual situation.

Storms are the clearest case on this creek. Rain changes conductivity, stage
and temperature together and quickly, rain-mode behaviour is unlabelled at
three of the four nodes, and most of the corpus is dry. Uber met the same
problem in traffic forecasting, where the departures that matter most, such as
holidays, are exactly the periods with the least history behind them, and
built a Bayesian network's uncertainty into anomaly detection for that
reason~\citep{zhu2017uber}. Storms play the part of holidays here.

It helps to name the two kinds of uncertainty involved~\citep{kendall2017uncertainties}.
\emph{Aleatoric} uncertainty is noise in the data itself, the wobble a
conductivity probe shows even when nothing is happening. More data does not
reduce it. \emph{Epistemic} uncertainty is uncertainty about the model, and it
is large where training data was thin and small where it was plentiful. More
data does reduce it. CobbleShoal already estimates the first through its
heteroscedastic scale head. Nothing in the repository estimated the second
before FlameSkimmer.

\subsection{The trained configuration}
\label{sec:configuration}

The values in Table~\ref{tab:config} are the settings of the first trained
checkpoint, read from the model definition, the training script and the
calibration artifact written beside the weights. The model is experimental and
not served; these are the settings the evaluation in
Section~\ref{sec:evaluation} measured.

\begin{table}[t]
  \centering
  \caption{The trained configuration of the first FlameSkimmer checkpoint.
    Values above the rule are shared with CobbleShoal's clock model; values
    below it belong to the Bayesian head and its training.}
  \label{tab:config}
  \begin{tabular}{ll}
    \toprule
    Setting & Value \\
    \midrule
    Sites $S$, canonical variables $V$, nodes $N$ & 5, 5, 17 \\
    Window $T$ & 24 timesteps (6 hours at 15 min sampling) \\
    Context features & 2 (hour sine and cosine) \\
    Hidden dimension $d$ & 64 \\
    Seed & 20260806 \\
    \midrule
    Head layers & $128 \to 32 \to 32 \to 2$, tanh between \\
    Prior & $\mathcal{N}(0, 1)$ on every head weight and bias \\
    Initial $\rho$ & $-5$, so $\sigma \approx 0.0067$ \\
    Supervised training targets $N_{\mathrm{train}}$ & 141{,}534 \\
    KL warmup & linear over the first 10\% of steps (752 steps) \\
    Head samples at prediction $S_{\mathrm{w}}$ & 64 \\
    Trainable parameters & 72{,}637, of which 10{,}500 in the head \\
    Operating false alarm rate $q$ & $10^{-3}$ \\
    Calibrated threshold $z_q$ & 16.271 \\
    \bottomrule
  \end{tabular}
\end{table}

The head is tiny next to the encoder, and that is on purpose. Its 5{,}250
weights and biases each carry a mean and a width, which is where the 10{,}500
comes from, and everything uncertain about the model lives in those numbers.
The 62{,}137 parameters of the encoder are ordinary point estimates.

\section{Data and Notation}
\label{sec:data}

\subsection{Corpus and split}
\label{sec:corpus}

FlameSkimmer trains on exactly the windows CobbleShoal's clock model trained on.
The corpus runs from 4 June 2025 to 29 April 2026, the span over which every
site in the roster reports. The split is chronological, with validation
starting on 1 February 2026, because a detector is deployed forward in time and
adjacent 15 minute windows overlap by 23 of their 24 steps, so a shuffled split
would put the training set into validation. Out of 19{,}096 fault-free anchors,
14{,}106 sit wholly on the training side and 4{,}990 on the validation side;
the run used evenly spaced subsets of 12{,}000 and 3{,}000. The recorded spans
are 4 June to 31 December 2025 for training and 25 February to 28 April 2026
for validation. The scaler is fitted on training rows only, and the two
footbridge sensors the training side never observed, dissolved oxygen and
floating conductivity, are held offline.

The corpus, the split, the anchors and the calibration routine are all imported
from CobbleShoal's training script rather than copied. A change there reaches
both models, and a comparison between them cannot be undone by the two copies
drifting apart.

\begin{table}[t]
  \centering
  \caption{Notation. Symbols shared with the CobbleShoal report keep their
    meaning there.}
  \label{tab:notation}
  \begin{tabular}{lll}
    \toprule
    Symbol & Meaning & Dimension \\
    \midrule
    $\bm{h}$ & encoder output for one node, node state and site context & $2d = 128$ \\
    $a$ & last reading of the node, the anchor & scalar \\
    $w$ & one head weight or bias & scalar \\
    $\mu_{w}, \rho_{w}$ & posterior mean and width parameter of $w$ & scalar \\
    $\sigma_{w}$ & posterior standard deviation, $\log(1 + e^{\rho_{w}})$ & scalar \\
    $s_{p}$ & prior standard deviation & scalar, 1 \\
    $q_{\phi}(w)$ & approximate posterior over head weights & \\
    $\beta$ & KL weight during warmup & scalar, $0 \to 1$ \\
    $\mu_{s}, \sigma_{s}$ & forecast and scale under head sample $s$ & $N$ \\
    $\bar{\mu}$ & predictive mean & $N$ \\
    $v_{\mathrm{al}}, v_{\mathrm{ep}}$ & aleatoric and epistemic variance & $N$ \\
    $v_{\mathrm{tot}}$ & total predictive variance & $N$ \\
    $\Lambda$ & combined Fisher statistic & $N$ \\
    \bottomrule
  \end{tabular}
\end{table}

\section{Model Architecture}
\label{sec:architecture}

\subsection{The shared trunk}
\label{sec:trunk}

The encoder is CobbleShoal's NestedEncoder, unchanged: decay LSTM over each
sensor node, sibling exchange within a site, directed transport downstream, and
a final mix back into the nodes. Its mathematics is set out in the CobbleShoal
report and is not repeated here.

What matters for FlameSkimmer is how it is built. The global random seed is set,
the encoder is constructed, and then the seed is set again with a fixed offset
of 10{,}000 before the head is constructed. CobbleShoal does the same. At the
same seed the two encoders therefore start from byte-identical weights, which a
unit test checks, and any difference between the two models' results belongs to
their heads rather than to a luckier draw of the trunk. FlameSkimmer composes
the encoder rather than inheriting CobbleShoal's base class, so none of the
leave-one-out or reconstruction machinery comes with it.

\subsection{The Bayesian head}
\label{sec:head}

The head reads the encoder output $\bm{h}$ for each node and produces two
numbers, a forecast delta and a raw scale. It has three layers,
\begin{equation}
  \bm{z}_{1} = \tanh(W_{1}\bm{h} + \bm{b}_{1}), \qquad
  \bm{z}_{2} = \tanh(W_{2}\bm{z}_{1} + \bm{b}_{2}), \qquad
  (\Delta, r) = W_{3}\bm{z}_{2} + \bm{b}_{3},
\end{equation}
and every entry of every $W$ and $\bm{b}$ is a random variable rather than a
number. The approximate posterior is mean-field Gaussian, so each weight is
independent of the others,
\begin{equation}
  q_{\phi}(w) = \mathcal{N}\big(\mu_{w},\ \sigma_{w}^{2}\big), \qquad
  \sigma_{w} = \log\big(1 + e^{\rho_{w}}\big),
\end{equation}
and the trainable parameters $\phi$ are the pairs $(\mu_{w}, \rho_{w})$. The
softplus keeps $\sigma_{w}$ positive without a constraint, which is the
parameterisation of Bayes by Backprop~\citep{blundell2015weight} and of the
bayesian-torch layers this head was modelled on~\citep{krishnan2022bayesiantorch}.
The means start where an ordinary linear layer's weights start, uniform within
$\pm 1/\sqrt{\text{fan in}}$, so the head begins at the same place
CobbleShoal's scale head does. Every $\rho_{w}$ starts at $-5$.

The forecast and its scale come out exactly as they do in CobbleShoal,
\begin{equation}
  \mu = a + \Delta, \qquad \sigma = \operatorname{softplus}(r) + 10^{-2},
\end{equation}
with the forecast anchored on the node's last reading and the scale floored so
the likelihood cannot reward claiming perfect certainty.

\subsection{Why only the head is Bayesian}
\label{sec:partial}

Making the encoder Bayesian too would multiply the parameters that need a
width, slow training, and force the whole encoder to be run once per sample at
prediction time. The evidence says it would not buy much. Across four inference
methods and eight datasets, networks with only some layers stochastic matched
fully stochastic ones, with no systematic benefit from going further~\citep{sharma2023stochastic}.
Keeping the trunk deterministic means the encoder runs once per window and only
the small head is sampled, so 64 samples cost 64 passes through about five
thousand weights rather than 64 passes through the whole model.

\subsection{Why two hidden layers}
\label{sec:depth}

A single hidden layer would be cheaper, and it would be wrong in a way that
matters here. Mean-field networks with one hidden layer provably cannot raise
their uncertainty in the gap between two well separated regions of
data~\citep{foong2020expressiveness}. Dry season and storm are, loosely, two
such regions, and the point of the model is the uncertainty between and beyond
them. The same work sees similar trouble empirically in deeper networks, which
is why the deep ensemble reference in Section~\ref{sec:limitations} exists.

\subsection{Keeping the heteroscedastic output}
\label{sec:hetero}

The head still predicts a scale for every reading. Replacing the per-reading
scale with weight uncertainty alone would lose the aleatoric part, and the
dispersion channel in Section~\ref{sec:channels} needs exactly that part.
Combining an input-dependent noise scale with weight uncertainty in one model is
the construction of~\citet{kendall2017uncertainties}. This is also why two
off-the-shelf alternatives were set aside: post hoc Laplace approximation and
variational Bayesian last layers, as implemented in their libraries, both learn
a single noise level shared by every input, which throws away the per-reading
scale.

\section{Training}
\label{sec:training}

\subsection{The evidence lower bound}
\label{sec:elbo}

Training fits $\phi$ by maximising the evidence lower bound,
\begin{equation}
  \mathcal{L}_{\mathrm{ELBO}}(\phi) =
    \mathbb{E}_{q_{\phi}(w)}\big[\log p(\mathcal{D} \mid w)\big]
    - \mathrm{KL}\big(q_{\phi}(w)\ \|\ p(w)\big),
\end{equation}
which trades fitting the data against staying close to the prior. The first
term rewards weights that explain the readings. The second charges for every
bit of certainty the posterior claims beyond what the prior allows, so a weight
the data says nothing about drifts back toward $\mathcal{N}(0, s_p^2)$, and its
width grows, which is where epistemic uncertainty comes from.

The likelihood is CobbleShoal's masked Gaussian negative log likelihood,
\begin{equation}
  \mathcal{L}_{\mathrm{nll}} = \frac{1}{|M|} \sum_{i \in M}
    \left[ \tfrac{1}{2} \left( \frac{y_i - \mu_i}{\sigma_i} \right)^{2} + \log \sigma_i \right],
\end{equation}
averaged over the entries $M$ where a target was observed and the anchor was
fresh enough to predict a delta from.

\subsection{The KL term, summed}
\label{sec:kl}

With a Gaussian posterior and a Gaussian prior the divergence has a closed form
per weight,
\begin{equation}
  \mathrm{KL}\big(\mathcal{N}(\mu_{w}, \sigma_{w}^{2})\ \|\ \mathcal{N}(0, s_{p}^{2})\big)
  = \log\frac{s_{p}}{\sigma_{w}} + \frac{\sigma_{w}^{2} + \mu_{w}^{2}}{2 s_{p}^{2}} - \frac{1}{2},
\end{equation}
and the KL of the whole head is the sum of this over all 5{,}250 weights and
biases. The sum matters. The bayesian-torch base layer returns the mean over
weights instead, which shrinks the penalty by a factor of the parameter count
and quietly turns the ELBO into something else. The same library also fills its
prior width with an argument named \texttt{prior\_variance}, so what it calls a
variance is used as a standard deviation. FlameSkimmer writes its own
small layer instead of depending on that library, which avoids both, and
avoids one more import paying the network file system cost on the OCF host. A
unit test checks the summed KL against PyTorch's own
\texttt{kl\_divergence}, and another checks that it comes out at zero when the
posterior equals the prior.

\subsection{Scaling and warmup}
\label{sec:beta}

The likelihood above is a mean over targets, so the KL has to be put on the same
per-target footing. The loss for one minibatch is
\begin{equation}
  \mathcal{L} = \mathcal{L}_{\mathrm{nll}} + \beta \, \frac{\mathrm{KL}(q_{\phi} \,\|\, p)}{N_{\mathrm{train}}},
\end{equation}
where $N_{\mathrm{train}} = 141{,}534$ is the number of supervised
$(\text{anchor}, \text{node})$ targets in the training set, counted with the
same mask the likelihood uses. Every minibatch carries the same share of the
KL. Blundell's paper also proposes a schedule that weights the KL by
$2^{M-i}/(2^{M}-1)$ for batch $i$ of $M$, and the BLiTZ library implements
it~\citep{esposito2020blitz}, but that schedule puts nearly all of the KL on
the first batch of each epoch and almost none on the rest, so it is not used.

The weight $\beta$ climbs linearly from 0 to 1 over the first 10\% of
optimiser steps and then stays at 1. Starting at zero lets the head learn a
sensible forecast before the prior starts pulling on it, and holding at 1
afterwards means the trained model optimises the true ELBO, with no tempering.

\subsection{Local reparameterization}
\label{sec:localrep}

The expectation in the ELBO is estimated by sampling. The obvious way is to draw
one weight matrix and push the whole minibatch through it, but then every
window in the batch shares the same noise, and the gradient estimate does not
get better as the batch grows. The local reparameterization
trick~\citep{kingma2015variational} samples the layer's output instead. For an
input $\bm{x}$ to a layer with independent Gaussian weights, each
pre-activation is itself Gaussian,
\begin{equation}
  b_{j} \sim \mathcal{N}\!\left(\sum_{i} x_{i}\,\mu_{ji} + \mu_{c_j},\ \
    \sum_{i} x_{i}^{2}\,\sigma_{ji}^{2} + \sigma_{c_j}^{2}\right),
\end{equation}
so one standard normal draw per output element gives an exact sample, with
independent noise for every window. The gradient variance then falls with batch
size as it should. The training forward uses this at every layer.

\subsection{Optimiser and stopping}
\label{sec:optim}

Training uses Adam at a learning rate of $10^{-3}$, batches of 64 anchors,
gradient clipping at 1.0, and CobbleShoal's stopping rule: the learning rate
halves after two epochs without improvement, and training stops after six,
keeping the best epoch. The quantity watched is the validation negative ELBO at
$\beta = 1$, with a fixed noise generator so the change between epochs is the
weights and not the draw. Table~\ref{tab:history} shows the run.

\begin{table}[t]
  \centering
  \caption{Training history. NLL is the training likelihood term, KL$/N$ the
    scaled divergence, and the last two columns the mean epistemic share on 600
    training and 600 validation anchors. Epoch 5 was kept.}
  \label{tab:history}
  \begin{tabular}{rrrrrrr}
    \toprule
    Epoch & NLL & KL$/N$ & $\beta$ & Val NLL & Share train & Share val \\
    \midrule
     1 & $-1.699$ & 0.166 & 0.25 & $-1.933$ & 0.160 & 0.147 \\
     2 & $-2.142$ & 0.162 & 0.50 & $-2.110$ & 0.095 & 0.080 \\
     3 & $-2.332$ & 0.159 & 0.75 & $-2.264$ & 0.155 & 0.126 \\
     4 & $-2.215$ & 0.156 & 1.00 & $-2.304$ & 0.131 & 0.122 \\
     5 & $-2.333$ & 0.152 & 1.00 & $-2.315$ & 0.135 & 0.115 \\
     6 & $-2.402$ & 0.148 & 1.00 & $-2.173$ & 0.076 & 0.058 \\
     7 & $-2.431$ & 0.145 & 1.00 & $-2.254$ & 0.123 & 0.091 \\
     8 & $-2.330$ & 0.142 & 1.00 & $-2.182$ & 0.124 & 0.097 \\
     9 & $-2.548$ & 0.140 & 1.00 & $-2.298$ & 0.101 & 0.071 \\
    10 & $-2.652$ & 0.139 & 1.00 & $-2.205$ & 0.129 & 0.099 \\
    11 & $-2.545$ & 0.137 & 1.00 & $-2.182$ & 0.108 & 0.078 \\
    \bottomrule
  \end{tabular}
\end{table}

The KL term never came close to dominating the loss, which was the first thing
to rule out. Its scaled size of about 0.15 sits beside a likelihood of about
$-2.3$.

\section{Prediction and Detection}
\label{sec:detection}

\subsection{The predictive distribution}
\label{sec:predictive}

At prediction time the encoder runs once and the head runs once for each of
$S_{\mathrm{w}} = 64$ sampled weight sets, giving a forecast $\mu_{s}$ and a
scale $\sigma_{s}$ per sample. The predictive distribution is a mixture of 64
Gaussians, and it is summarised by its first two moments,
\begin{align}
  \bar{\mu} &= \frac{1}{S_{\mathrm{w}}} \sum_{s} \mu_{s}, \\
  v_{\mathrm{al}} &= \frac{1}{S_{\mathrm{w}}} \sum_{s} \sigma_{s}^{2}, \qquad
  v_{\mathrm{ep}} = \frac{1}{S_{\mathrm{w}}} \sum_{s} \big(\mu_{s} - \bar{\mu}\big)^{2}, \qquad
  v_{\mathrm{tot}} = v_{\mathrm{al}} + v_{\mathrm{ep}} .
\end{align}
This is the law of total variance. The aleatoric part is the noise each version
of the model expects, averaged; the epistemic part is how much the versions
disagree about the forecast itself. The variance is the population form,
dividing by $S_{\mathrm{w}}$, because that is what the moment-matched mixture
uses. The ratio $v_{\mathrm{ep}} / v_{\mathrm{tot}}$, the \emph{epistemic share},
is the number the evaluation watches.

\subsection{Sampling whole weights at prediction time}
\label{sec:weightspace}

Training samples outputs; prediction samples whole weight matrices, and the
difference is deliberate. Local reparameterization gives every window in a
batch its own independent noise, which is what training wants and what scoring
cannot have: a window's score would then depend on which other windows happened
to share its batch and where it sat in it. Drawing $S_{\mathrm{w}}$ complete
weight sets and passing every window through the same ones removes that. For
any single window the two schemes give exactly the same distribution, since
each layer's weights are used once per window and are independent of the layers
before them. A unit test checks that a window scores identically alone and
stacked beside another.

\subsection{Fixed serving noise}
\label{sec:crn}

The 64 weight sets are drawn once, when the model is constructed, from a fixed
seed and a private generator, and reused on every call. Two runs over the same
window therefore give the same score to the last bit, which matters because an
alert that appears on one run and vanishes on the next is worse than either
answer. Drawing from a private generator leaves the global random state
untouched, so constructing the model does not disturb anything seeded after it.
What is stored is the standard normal noise rather than the weights themselves,
so the same draw stays valid however training later moves $\mu$ and $\rho$.

\subsection{Two channels, and which variance each reads}
\label{sec:channels}

FlameSkimmer has the two channels that come free with a forecast head, and the
rule that matters most in the whole design is which variance each one reads.

\textbf{Excursion} is the standardised one-step residual against the
\emph{total} predictive standard deviation,
\begin{equation}
  c_{\mathrm{exc}} = \frac{|y - \bar{\mu}|}{\sqrt{v_{\mathrm{tot}}}} .
\end{equation}
A reading far from the forecast in conditions the model knows well should count
for a lot. The same distance in conditions the model admits it does not know
should count for less, and dividing by the total variance does exactly that.

\textbf{Dispersion} reads the \emph{aleatoric} part only,
\begin{equation}
  c_{\mathrm{disp}} = \max\!\left(0,\ \log \frac{\max(\sqrt{v_{\mathrm{al}}}, \varphi)}{\max(\hat{s}, \varphi)}\right),
\end{equation}
where $\hat{s}$ is the observed spread over the last 8 steps and
$\varphi = 10^{-3}$ is a floor. Dispersion asks whether a sensor is flatter than
its own noise, which is how a frozen probe gives itself away. If the epistemic
variance were included, the expected spread would swell whenever the model was
unsure, an ordinary sensor would then look too flat by comparison, and every
storm would read as a stuck sensor. The functions computing both channels are
CobbleShoal's frozen ones, called with different scales; their arithmetic is
not touched, because the calibrated artifacts depend on it.

\subsection{Nulls, combination and threshold}
\label{sec:calibration}

Everything after the channels is CobbleShoal's procedure, run on FlameSkimmer's
scores. Each channel becomes a right-tail probability against its own
fault-free distribution, empirical below the 95th percentile and a fitted
Generalised Pareto above it~\citep{pickands1975gpd}. The two probabilities are
combined per node as $\Lambda = -2 \sum_{j} \log p_{j}$~\citep{fisher1925methods},
and the threshold comes from peaks over threshold on the fault-free values of
$\Lambda$~\citep{siffer2017spot},
\begin{equation}
  z_{q} = u - \sigma \log\!\left(\frac{q\,n}{N_{u}}\right) \quad (\gamma = 0).
\end{equation}
The fit used 400 fault-free training windows, giving $n = 4{,}769$ node scores,
an initial threshold $u = 10.135$ with $N_{u} = 96$ exceedances,
$\sigma = 2.044$ and $\gamma = 0$, so $z_{q} = 16.271$ at $q = 10^{-3}$.
Measured on 2{,}118 held-out validation scores the exceedance rate was
$4.7 \times 10^{-4}$.

Rain is handled outside the model, as it is for CobbleShoal. The rain gate moves
the threshold and never the score, and FlameSkimmer declares no built-in support
behaviour. The comparison in the next section uses the ungated threshold for
both models.

\section{Evaluation}
\label{sec:evaluation}

FlameSkimmer was compared with CobbleShoal's clock model, trained by the same
script on the same split from the same seed. Every number is in the scaler's
standardised units, and the evaluation checks that both models were scaled
identically before comparing them. Results are reported per node; a mean across
nodes hides the nodes where a model is worse.

\subsection{Forecast quality}
\label{sec:quality}

Table~\ref{tab:quality} gives the negative log likelihood, the continuous
ranked probability score and the coverage of the central 90\% interval on the
3{,}000 validation anchors. The CRPS uses the closed form for a
Gaussian~\citep{gneiting2007scoring}, applied to the moment-matched predictive
distribution, and was copied from uncertainty-toolbox~\citep{chung2021toolbox}
rather than adding the package. The NLL here includes the
$\tfrac{1}{2}\log 2\pi$ constant the training loss leaves out and is taken over
the full predictive distribution, so it is not directly comparable with the
validation NLL in Table~\ref{tab:history}.

\begin{table}[t]
  \centering
  \caption{Forecast quality on validation, FlameSkimmer against CobbleShoal.
    Lower NLL and CRPS are better; the 90\% interval should cover close to
    90\%. Footbridge has only 21 supervised targets on the validation side.}
  \label{tab:quality}
  \begin{tabular}{lrrrrrrr}
    \toprule
     & & \multicolumn{2}{c}{NLL} & \multicolumn{2}{c}{CRPS} & \multicolumn{2}{c}{Coverage 90} \\
    Node & $n$ & Flame & Cobble & Flame & Cobble & Flame & Cobble \\
    \midrule
    all nodes                 & 31{,}650 & $-1.460$ & $-1.157$ & 0.0335 & 0.0330 & 0.972 & 0.961 \\
    \midrule
    north\_fork\_0 cond.      & 2{,}996 & $-1.588$ & $-1.410$ & 0.0243 & 0.0277 & 0.972 & 0.977 \\
    north\_fork\_0 depth      & 2{,}996 & $-0.503$ &  $0.461$ & 0.0470 & 0.0522 & 0.977 & 0.988 \\
    north\_fork\_0 temp.      & 2{,}996 & $-1.432$ &  $0.275$ & 0.0178 & 0.0154 & 0.983 & 0.894 \\
    footbridge cond.          &      21 & $-1.222$ &  $0.045$ & 0.0383 & 0.1041 & 1.000 & 1.000 \\
    footbridge depth          &      21 & $-1.052$ & $-0.590$ & 0.0490 & 0.0650 & 0.952 & 1.000 \\
    footbridge temp.          &      21 & $-1.712$ & $-1.443$ & 0.0220 & 0.0252 & 1.000 & 1.000 \\
    south\_fork\_1 cond.      & 1{,}536 & $-1.382$ & $-0.681$ & 0.0330 & 0.0543 & 0.958 & 0.980 \\
    south\_fork\_1 depth      & 1{,}536 & $-0.540$ & $-0.530$ & 0.0701 & 0.0688 & 0.993 & 0.991 \\
    south\_fork\_1 temp.      & 1{,}536 & $-1.645$ & $-1.755$ & 0.0235 & 0.0206 & 0.990 & 0.978 \\
    south\_fork\_2 cond.      & 2{,}997 & $-1.924$ & $-1.664$ & 0.0217 & 0.0258 & 0.973 & 0.939 \\
    south\_fork\_2 depth      & 2{,}997 & $-0.903$ & $-1.030$ & 0.0733 & 0.0496 & 0.989 & 0.985 \\
    south\_fork\_2 temp.      & 2{,}997 & $-2.115$ & $-2.336$ & 0.0155 & 0.0125 & 0.960 & 0.934 \\
    oxford cond.              & 3{,}000 & $-1.980$ & $-1.894$ & 0.0157 & 0.0166 & 0.984 & 0.985 \\
    oxford depth              & 3{,}000 & $-0.972$ & $-0.899$ & 0.0570 & 0.0590 & 0.907 & 0.968 \\
    oxford temp.              & 3{,}000 & $-2.141$ & $-2.184$ & 0.0160 & 0.0148 & 0.985 & 0.943 \\
    \bottomrule
  \end{tabular}
\end{table}

The picture is mixed. FlameSkimmer is better on likelihood overall and better
on both NLL and CRPS at every conductivity node, which is the variable that is
scored for detection. It is worse on NLL at four nodes and on CRPS at six,
mostly temperature and stage on the south fork and at oxford, so the proposed
bar of no node worse on either measure is not met. Its intervals are also too
wide. The central 90\% interval covers 97.2\% of validation targets, outside
the 85 to 95\% band that would count as honest. Wide intervals improve the
likelihood on rough nodes and cost CRPS on smooth ones, which fits the pattern
in the table.

\subsection{Epistemic behaviour}
\label{sec:epistemic}

The test the model exists to pass is whether its epistemic share rises where it
has seen little. A step counts as rain when more than 0.1 mm fell over the
preceding six hours, from the Open-Meteo archive. Table~\ref{tab:epistemic}
gives the share on every supervised target, split by side and by rain, along
with the median epistemic standard deviation in absolute terms.

\begin{table}[t]
  \centering
  \caption{Epistemic share $v_{\mathrm{ep}} / v_{\mathrm{tot}}$ per supervised
    target, and the median epistemic standard deviation
    $\sqrt{v_{\mathrm{ep}}}$. The three event rows count rain steps only; the
    January event is the botanical garden actuator.}
  \label{tab:epistemic}
  \begin{tabular}{lrrrr}
    \toprule
    Group & $n$ & Median share & Quartiles & Median $\sqrt{v_{\mathrm{ep}}}$ \\
    \midrule
    Training, dry      & 131{,}459 & 0.132 & 0.082 to 0.215 & 0.0149 \\
    Training, rain     &  10{,}075 & 0.066 & 0.008 to 0.151 & 0.0147 \\
    Validation, dry    &  29{,}385 & 0.132 & 0.049 to 0.184 & 0.0149 \\
    Validation, rain   &   2{,}265 & 0.027 & 0.002 to 0.117 & 0.0147 \\
    \midrule
    Event, November 2025 rain & 813 & 0.021 & 0.001 to 0.084 & \\
    Event, April 2026 rainfall & 810 & 0.064 & 0.003 to 0.157 & \\
    Event, January 2026 actuator & 5{,}664 & 0.022 & 0.002 to 0.092 & \\
    \bottomrule
  \end{tabular}
\end{table}

The share falls in rain rather than rising, on both sides of the split and in
all three event windows. The last column says why. The absolute epistemic
standard deviation is 0.0149 or 0.0147 in every group, so the epistemic
variance is effectively a constant. What moves in rain is the aleatoric part:
the scale head correctly predicts a noisier sensor, the total grows, and a
constant divided by a larger total gets smaller. The model is not saying it is
more certain in storms. It is not saying anything with its epistemic variance
at all.

\subsection{Detection}
\label{sec:detectionresults}

Table~\ref{tab:events} scores every labelled event in the catalogue. For each
entry the peak combined score at the named site's conductivity node, over every
anchor in the event window, is compared with each model's threshold at
$q = 10^{-3}$: 16.271 for FlameSkimmer and 29.070 for CobbleShoal. The two
thresholds are on different scales, since one combines two channels and the
other four.

\begin{table}[t]
  \centering
  \caption{Peak score at the graded node in each labelled event window. Entries
    labelled relative only are reported, not graded.}
  \label{tab:events}
  \begin{tabular}{llrrl}
    \toprule
    Event & Site & Flame & Cobble & Label \\
    \midrule
    June 2025 discharge         & south\_fork\_1 & 17.655 & 34.539 & anomaly, both flagged \\
    September 2025 overnight    & south\_fork\_1 & 23.740 & 47.566 & anomaly, both flagged \\
    September 2025 overnight    & south\_fork\_2 & 15.796 & 28.405 & anomaly, both missed \\
    March 2026 hydrant          & north\_fork\_0 & 19.467 & 46.026 & anomaly, both flagged \\
    \midrule
    June 2025 discharge         & south\_fork\_2 & 17.880 & 26.872 & relative only \\
    August 2025 sprinklers      & north\_fork\_0 & 11.228 & 13.507 & relative only \\
    November 2025 foam          & north\_fork\_0 & 23.387 & 48.388 & relative only \\
    \bottomrule
  \end{tabular}
\end{table}

Both models flag three of the four labelled anomalies and both miss the
south\_fork\_2 half of the September event, FlameSkimmer by about 3\% of its
threshold. The margins are thinner for FlameSkimmer, though. Its peaks sit 8 to
46\% above its threshold, where CobbleShoal's sit 19 to 64\% above its own,
which is what dropping the leave-one-out and reconstruction channels would be
expected to cost.

\subsection{False alarm rate}
\label{sec:far}

Table~\ref{tab:far} gives the rate at which fault-free validation node scores
crossed each model's threshold at $q = 10^{-3}$, overall and split by rain.

\begin{table}[t]
  \centering
  \caption{Empirical false alarm rate on fault-free validation anchors at
    $q = 10^{-3}$. CobbleShoal has more finite scores because its leave-one-out
    channel does not need a fresh anchor.}
  \label{tab:far}
  \begin{tabular}{lrrrrrr}
    \toprule
     & \multicolumn{2}{c}{All} & \multicolumn{2}{c}{Dry} & \multicolumn{2}{c}{Rain} \\
    Model & Rate & $n$ & Rate & $n$ & Rate & $n$ \\
    \midrule
    FlameSkimmer & $1.07 \times 10^{-3}$ & 31{,}722 & $7.5 \times 10^{-4}$ & 29{,}445 & $5.3 \times 10^{-3}$ & 2{,}277 \\
    CobbleShoal  & $1.36 \times 10^{-3}$ & 37{,}626 & $6.9 \times 10^{-4}$ & 34{,}905 & $9.9 \times 10^{-3}$ & 2{,}721 \\
    \bottomrule
  \end{tabular}
\end{table}

The overall rate is close to nominal for both models, well inside the order of
magnitude that the peaks over threshold diagnostics already treat as
acceptable. In rain FlameSkimmer flags at about half CobbleShoal's rate. That
looks like the payoff the model was built for, but Section~\ref{sec:epistemic}
rules out the intended explanation, since the epistemic variance does not change
in rain. The likelier cause is the wider total intervals of
Section~\ref{sec:quality}, which damp the excursion channel everywhere and
matter most where readings move most.

\subsection{The gate, and why it failed}
\label{sec:gate}

The plan set one test for whether the Bayesian head is doing its job: the
median epistemic share on validation rain steps should be at least 1.5 times
the median on training dry steps. It is $0.027 / 0.132 = 0.20$ times.

The cause is in the trained widths. Table~\ref{tab:sigma} gives the posterior
standard deviations after training. Every one of them started at
$\log(1 + e^{-5}) \approx 0.0067$, and they have barely moved.

\begin{table}[t]
  \centering
  \caption{Posterior standard deviations of the head after training. Every value
    started at 0.0067.}
  \label{tab:sigma}
  \begin{tabular}{lrrr}
    \toprule
    Layer & Median $\sigma$, weights & Largest $\sigma$, weights & Median $\sigma$, biases \\
    \midrule
    $128 \to 32$ & 0.0129 & 0.0160 & 0.0093 \\
    $32 \to 32$  & 0.0064 & 0.0108 & 0.0061 \\
    $32 \to 2$   & 0.0053 & 0.0068 & 0.0036 \\
    \bottomrule
  \end{tabular}
\end{table}

The ceiling on the first layer is not a coincidence. Adam moves each parameter
by roughly the learning rate per step at most, and the kept epoch came after
about 940 steps at a learning rate of $10^{-3}$. Starting from $-5$, no
$\rho_{w}$ could have climbed much past $-4.06$, which is
$\sigma_{w} \approx 0.017$. The largest width actually learned is 0.016. The
widths were limited by the optimiser and the starting point, not chosen by the
data, and early stopping kept the epoch immediately after $\beta$ reached 1,
before the prior had any time at full strength.

That also explains why the epistemic variance is constant. When every width
sits near the same small value, the variance the head adds is roughly that
width squared times the squared size of the activations feeding each layer, and
after a tanh those activations have much the same size for almost every input.
The variance that results barely depends on the window, so it cannot rise in a
storm.

None of this says that a Bayesian head cannot do the job. It says this one was
not trained long enough, or started too narrow, for its widths to learn
anything.

\section{Limitations and Future Work}
\label{sec:limitations}

The first thing to fix is the training of the widths. Three changes address
Section~\ref{sec:gate} directly: start $\rho$ less negative, for example at
$-3$, which gives $\sigma \approx 0.049$; give the $\rho$ parameters their own,
larger learning rate; and stop early only after $\beta$ has been at 1 for
several epochs. Each is a change to a default rather than to the design, and
the gate should be measured again after them before anything else is built on
the model.

The plan also calls for a reference when the gate fails: a five-seed deep
ensemble of the same trunk with CobbleShoal's deterministic scale head, reported
beside FlameSkimmer. Ensembles match or beat approximate Bayesian networks in
most comparisons~\citep{lakshminarayanan2017ensembles} and held up best under
dataset shift in the largest benchmark of the question~\citep{ovadia2019shift},
though that benchmark is classification. The ensemble is not Bayesian over
weights, but it is the bar a Bayesian head should clear, and it has not yet
been run. Monte Carlo dropout, the other cheap option, was the weakest of four
methods in a rainfall-runoff LSTM benchmark~\citep{klotz2022rainfall} and is
not planned.

The chronological split was expected to put most of the rain on the validation
side. By anchor count it does not: 7.0\% of training anchors and 7.3\% of
validation anchors had rain in the preceding six hours. The storms may still
differ in size or season between the two sides, but the split is a weaker test
of shift than it was taken to be, and a blocked split or a held-out storm would
make a sharper one.

Only one seed has been trained. Every result in Section~\ref{sec:evaluation}
should be repeated across several before it is relied on, and reported per seed.

The event catalogue has no true negatives, no storm that the group has labelled
as a known non-event. The rain false alarm comparison therefore rests on
fault-free statistics alone. One labelled storm would turn it into a real case.

FlameSkimmer is not served. A serving detector, its entry in the serving
registry and updates to the README and user manual follow only once the
evaluation is accepted. Whether the epistemic signal should ever become an
alert type of its own is a separate decision, and nothing in this design adds
an uncertainty rule to the alerting path.

Finally, every number in this document describes one experimental checkpoint
and the calibration fitted alongside it. They will move when the model is
retrained, and some of them, by design, should.

\bibliographystyle{unsrtnat}
\bibliography{references}

\end{document}
